\documentclass[10pt,a4paper]{amsart}
\usepackage[margin=19mm]{geometry}
\usepackage{amsmath,amssymb,mathtools}
\usepackage[scr=rsfs,cal=euler]{mathalfa}
\usepackage{xfrac}
\usepackage[unicode,hidelinks,bookmarks=false]{hyperref}
\newcommand{\ie}{i.e.\ }
\newcommand{\cf}{cf.\ }
\newcommand{\resp}{respectively\ }
\newcommand{\Subsection}{\subsection}
\providecommand{\nobreakdash}{\nobreak\hbox{-}}
\setlength{\emergencystretch}{2em}
\multlinegap0pt
\usepackage{bm}
\usepackage{bbm}
\usepackage{dsfont}
\usepackage{xspace}
\usepackage{cancel}
\usepackage{mathtools}
\usepackage{amsthm}
\makeatletter
\@ifundefined{theorem}{\newtheorem{theorem}{Theorem}}{}
\@ifundefined{lemma}{\newtheorem{lemma}[theorem]{Lemma}}{}
\@ifundefined{proposition}{\newtheorem{proposition}[theorem]{Proposition}}{}
\makeatother
\title[Ascent and descent of Artinian module structures under flat ]{Ascent and descent of Artinian module structures under flat base changes}
\author{Tran Do Minh Chau, Le Thanh Nhan}
\date{Source: arXiv:2507.19019v1 (CC BY 4.0)}
\begin{document}
\maketitle

\noindent\emph{Source and licence.} Ascent and descent of Artinian module structures under flat base changes. Version arXiv:2507.19019v1 (CC BY 4.0), distributed under the Creative Commons Attribution 4.0 International licence, as recorded in the arXiv licence metadata for this exact version and shown on the abstract page https://arxiv.org/abs/2507.19019v1. Full rights notice: licenses/arxiv-2507.19019-CC-BY-4.0.txt.\par\smallskip

\noindent\textit{Editorial scope.} Counted unit: the Theorem whose source label is \texttt{T:2} in the article (source0.tex lines 297--298, source label \texttt{T:2}), delivered together with the article's own complete original proof of it (source0.tex lines 300--322, one \texttt{proof} environment of 23 lines). No mathematics is written, repaired or re-derived: statement and proof are the article's own text line for line. Required context retained uncounted: L:1a (lines 164--177); P:1a (lines 189--200). Every definition, display and label of those lines is retained unchanged. Omitted from this package and not redistributed: 1--163, 178--188, 201--296, 299--299, 323--371. Nothing inside a retained excerpt is omitted or altered: every retained body is delivered exactly as the article's lines are written, so no omission receipt of any kind accompanies this package. Citations: the retained excerpts carry no citation command at all, so no bibliography entry of the article is transcribed and no reference is redistributed. Reference adaptation: none was needed and none was made; every reference inside the delivered unit resolves inside it, so no unresolved reference or citation is produced. Printed slips kept literal: no printed word, symbol, hypothesis or number of the counted statement or of its proof was altered. Layout and class: the article's own class is \texttt{article} with packages including amsmath, amsthm, amsxtra, amssymb, latexsym, amscd, eucal; the wrapper is the lane's pinned \texttt{amsart} editorial wrapper at 10pt on A4, which restates the theorem-like environments the retained text needs and adds the article's own macro definitions that the retained text needs (none), with extra wrapper packages (bm, bbm, dsfont, xspace, cancel, mathtools). The delivered numbering is the unit's own. Assets: no retained span contains a figure environment, a graphics-inclusion command, a PDF-page inclusion, a table or a TikZ picture; no image file is copied into this package, the package declares no asset dependency and no image file is redistributed with it. Licence: version arXiv:2507.19019v1 of this article is distributed under the Creative Commons Attribution 4.0 International licence, as recorded in the arXiv licence metadata for this exact version and shown on the abstract page https://arxiv.org/abs/2507.19019v1. A full rights notice is delivered at licenses/arxiv-2507.19019-CC-BY-4.0.txt. Characters: every retained excerpt is pure ASCII, so the delivered engine, XeLaTeX with its default Latin Modern text font, reports no missing character.
\begin{lemma} \label{L:1a} The following statements are equivalent:

(a)  Every  finite length $S$-module is of finite length as an $R$-module (by means of $\varphi$).

(b) Every  finite length $S$-module is Artinian as an $R$-module.

(c) There exists a finite length $S$-module $B\neq 0$  such that $B$ is Artinian as an $R$-module.

(d) $\ell_R (S/\frak n)<\infty$.


Suppose that the equivalent conditions (a), (b), (c), (d) satisfy. If $B$ is a finite length $S$-module, then 
$$\ell_R(B)=\ell_S(B)\ell_R(S/\frak n).$$
\end{lemma}

\begin{proposition} \label{P:1a} The following statements are equivalent:

(a)  $\dim(S/\frak mS)=0.$

(b)   $A\otimes_RS$ is an Artinian $S$-module for every Artinian $R$-module $A$. 

(c) $A\otimes_RS$ is an Artinian $S$-module for some Artinian $R$-module $A\neq 0$.

(d) $A\otimes_RS$ is a finite length $S$-module for every finite length $R$-module $A$.
 
(e) $A\otimes_RS$ is a finite length $S$-module for some finite length $R$-module $A\neq 0$.
\end{proposition}

\begin{theorem} \label{T:2} Assume that $\ell_R(S/\frak mS)=m$. Let $A$ be an Artinian $R$-module. Then  $A^m$ has a structure of Artinian $S$-module and there is an isomorphism $A\otimes_RS\cong A^m$ of $S$-modules. If we regard this $S$-module $A^m$ as an $R$-module by means of $\varphi$, then we recover the usual $R$-module structure on $A^m$.  
\end{theorem}

\begin{proof} Note that $S/\frak mS$ is a $R/\frak m$-vector space of dimension $m$.  Let $\epsilon_1, \ldots , \epsilon_m$ be elements in $S$ such that $\epsilon_1+\frak mS, \ldots , \epsilon_m+\frak mS$ is a  base of  the vector space $S/\frak mS$. We claim that $S/\frak m^tS=\big(R(\epsilon_1, \ldots , \epsilon_m)+\frak m^tS\big)/\frak m^tS$ for all $t\in\Bbb N$. Indeed, since $S=R(\epsilon_1, \ldots , \epsilon_m)+\frak mS,$ we have 
$$S/\frak m^tS=\big(R(\epsilon_1, \ldots , \epsilon_m)+\frak m^tS\big)/\frak m^tS + \frak m(S/\frak m^tS).$$ 
It follows by Lemma \ref{L:1a} and by the faithful flatness of $S$ over $R$ that 
\begin{align}\ell_R(S/\frak m^tS)=\ell_S(S/\frak m^tS)\ell_R(S/\frak n)&=\ell_R(R/\frak m^t)\ell_S(S/\frak mS)\ell_R(S/\frak n)\notag\\
&=\ell_R(R/\frak m^t)\ell_R(S/\frak mS)=m\ell_R(R/\frak m^t)=\ell_R\big((R/\frak m^t)^m\big)<\infty.\notag\end{align} Therefore, the claim follows by Nakayama Lemma. We consider  $S/\frak m^tS$ as an $R/\frak m^t$-module.  Then, by the claim we can define a surjection  $\varphi_t: (R/\frak m^t)^m\rightarrow S/\frak m^tS$ given by $$\varphi_t(r_1+\frak m^t, \ldots , r_m+\frak m^t)=(r_1\epsilon_1+ \ldots +r_m\epsilon_m)+\frak m^tS$$ for all $r_1, \ldots , r_m\in R.$ Since $\ell_R(S/\frak m^tS)=\ell_R\big((R/\frak m^t)^m\big),$ it follows that $\varphi_t$ is an isomorphism.


  Consider $A^m$ as an $R$-module with the usual scalar multiplication. Define a bilinear map $f: A\times S\rightarrow A^m$ as follows: Let $a\in A$ and $s\in S.$ Since $A$ is Artinian, there exists $t\in\Bbb N$ such that $\frak m^ta=0$. By the claim, there exist $r_1, \ldots , r_m\in R$ such that $s=\sum_{i=1}^m r_i\epsilon_i \ ({\rm mod} (\frak m^tS))$. Set $f(a, s):=(r_1a, \ldots , r_ma).$ As $\varphi_t$ is isomorphic,  we can check that the  definition of  $f(a, s)$ does not depend on the choice of  $t$ and $r_1, \ldots , r_m$.  Therefore, there exists an $R$-homomorphism $h: A \otimes_RS\rightarrow A^m$ given by $h(a\otimes s)=f(a, s)$ for all $a\in A$ and $s\in S.$ Let $(a_1, \ldots , a_m)\in A^m$, we have 
 $$(a_1, \ldots , a_m)=\sum_{i=1}^m\big(f(a_i,\epsilon_i)\big)=\sum_{i=1}^m h(a_i\otimes \epsilon_i)=h\big(\sum_{i=1}^m(a_i\otimes \epsilon_i)\big).$$ Therefore, $h$ is a surjection. 

We consider a $R$-homomorphism $g: A^m\rightarrow A\otimes_RS$  given by $g(a_1, \ldots , a_m):=\sum_{i=1}^m(a_i\otimes\epsilon_i)$ for any $(a_1, \ldots , a_m)\in A^m$. Let $a\in A, s\in S$ and let  $t\in\Bbb N$ be such that $\frak m^ta=0$.  By the claim, there exist $r_1, \ldots , r_m\in R$ such that $s=\sum_{i=1}^m r_i\epsilon_i\ ({\rm mod}(\frak m^tS)).$ Note that 
$a\otimes s=a\otimes \sum_{i=1}^m (r_i\epsilon_i)=\sum_{i=1}^m r_i(a\otimes \epsilon_i).$ Therefore,  $\{a\otimes \epsilon_i\mid a\in A, 1\leq i\leq m\}$ is a system of generators of $R$-module $A\otimes_R S$. By the definitions of $f, g, h$ we have 
$$gh(a\otimes \epsilon_1)=gf(a,\epsilon_1)=g(a,0, \ldots ,0)=a\otimes\epsilon_1.$$ Similarly, we have $gh(a\otimes \epsilon_i)=a\otimes\epsilon_i$ for all $2\leq i\leq m.$ It follows that $gh={\rm id}_{A\otimes_RS}.$ Therefore, $h$ is injective. Hence, $h$ is an isomorphism of $R$-modules and $hg={\rm id}_{A^m}.$ 

Now we provide a  $S$-module structure on $A^m$ as follows. Let $\underline a=(a_1, \ldots , a_m)\in A^m$ and $s\in S.$  Set $$s\circ\underline a:=h\big(s(g(\underline a))\big).$$
For any $s, s'\in S$, $\underline a, \underline a'\in A^m$, it is clear that $(s+s')\circ\underline a=s\circ\underline a+s'\circ\underline a$, $s\circ(\underline a+\underline a')=s\circ\underline a +s\circ\underline a'$. Since $hg={\rm id}_{A^m},$ we have $1\circ\underline a=h(1g(\underline a))=hg(\underline a)=\underline a.$ Moreover, 
$$s'\circ(s\circ\underline a)=h(s'g(s\circ\underline a))=h(s'g(h(sg(\underline a))))=h(s'sg(\underline a))=(s's)\circ\underline a.$$ So, this is a scalar multiplication and it makes $A^m$ being a $S$-module. Consider this $S$-module $A^m$ as an $R$-module by means of $\varphi$, then for $r\in R$ and $\underline a=(a_1, \ldots , a_m)\in A^m$ we have  
\begin{align}\varphi(r)\circ\underline a=h(\varphi(r)g(\underline a))&=h(\varphi(r)\sum_{i=1}^m(a_i\otimes\epsilon_i))=\sum_{i=1}^m h(a_i\otimes \varphi(r)\epsilon_i)\notag\\
&=\sum_{i=1}^m h(a_i\otimes r\epsilon_i)=\sum_{i=1}^m h(ra_i\otimes \epsilon_i)=(ra_1, \ldots , ra_m).\notag\end{align} Therefore, we recover the usual $R$-module structure on $A^m$. 

Let $s\in S$ and $z\in A\otimes_RS.$ Then $\displaystyle z=\sum_{i=1}^m u_i(a_i\otimes \epsilon_i)$ for some elements $u_i\in R$ and $a_i\in A$. By the definition of $h$,  $$h(z)=\sum_{i=1}^m h(u_ia_i\otimes \epsilon_i)=(u_1a_1, \ldots , u_ma_m).$$ Since $gh={\rm id}_{A\otimes_RS},$ we get by the definition of the scalar multiplication of $A^m$ over $S$ that 
 $$s\circ h(z)=s\circ (u_1a_1, \ldots , u_ma_m)=h(sg(u_1a_1, \ldots , u_ma_m))=h(sgh(z))=h(sz).$$ Therefore, $h$ is also an isomorphism of $S$-modules. Since $A\otimes_RS$ is an Artinian $S$-module by Proposition \ref{P:1a}, it follows that  $A^m$ is also an Artinian $S$-module.
\end{proof}

\end{document}
